\documentclass[11pt]{article}
\usepackage[margin=1in]{geometry}
\usepackage{amsmath,amssymb,parskip}
\usepackage[colorlinks=true,allcolors=blue]{hyperref}
\newcommand{\E}{\mathbb E}
\setlength{\parskip}{5pt}
\title{Earnings-Aware Option Pricing}
\author{}\date{}
\begin{document}
\maketitle
\vspace{-3.2em}

\section{The Main Idea}

Black Scholes often fails during earnings events and sometimes for overnight returns. 
This failure occurs because the model assumes a constant volatility, which can't really be
properly scaled for different events or earnings reports. The stock may not move much 
during the days leading up to the event, before gapping 5-10\% overnight. This model 
attempts to fix this by splitting the stock's variance into two parts. It prices the 
ordinary day-to-day movement separately from the single jump on the known earnings date.
Optionally, the diffusion part can also have stochastic volatility (Heston) to produce a realistic skew.

Notation: spot $S_0$, strike $K$, years to expiry $T$, rate $r$, dividend yield
$q$. Forward $F=S_0e^{(r-q)T}$, discount $D=e^{-rT}$. Earnings happens at a fixed
time $t_e$.

\section{The Earnings Jump}

Standard ways to process jumps are either Merton's jump-diffusion or Kou's double-exponential.
However, both of these are continuous-time models with a Poisson process for the jump. 
In this case, we replace this jump J with K different normals, one per "scenario" for how the earnings report lands.

\begin{equation}
f_J(x) = \sum_{i=1}^{K} w_i\,\mathcal N(x;\,\mu_i,\,s_i^2).
\end{equation}

This represents a weighted average of bell curves. Each scenario $i$ has a weight $w_i$ (the probability of that scenario, summing to 1), a mean $\mu_i$ (the expected jump size for that scenario), and a variance $s_i^2$ (the uncertainty around that jump size).
This allows you to represent different possible outcomes of the earnings report, such as a beat, miss, or in-line result, and their associated probabilities and uncertainties.


\textbf{Keeping it arbitrage-free.} Under the risk-neutral measure the jump can't
add any drift to the stock, which means it has to satisfy $\E[e^J]=1$. Given the
scenario centres $m_i$ you actually want, you shift all of them by the same
constant so this holds:
\begin{equation}
\mu_i = m_i - \log\!\Big(\textstyle\sum_i w_i\,e^{\,m_i+s_i^2/2}\Big).
\end{equation}
Shifting every centre by the same amount keeps the shape (the spacing and widths)
and just recenters the whole distribution. Note this is a pricing adjustment, not a claim
about real-world probabilities: a distribution fit to historical returns is not
automatically the risk-neutral one.

\section{Pricing with no simulation}

Usually, adding jumps means Monte Carlo simulations across many price paths and averaging
the payoffs. However, for this, that is not needed as the event date is fixed. 

Suppose scenario $i$ happened. Then over the life of the option the log-return is
the ordinary diffusion (a normal) plus that scenario's jump (also a normal). The
sum of two independent normals is normal, so on that branch the terminal price is
lognormal --- exactly the same as Black--Scholes. So each branch prices with the
standard Black-76 formula, and the option value is the probability-weighted sum
over branches:
\begin{equation}
C = D\sum_{i=1}^{K} w_i\,\big[F_i\,\Phi(d_{1i})-K\,\Phi(d_{2i})\big],
\end{equation}
where $\Phi$ is the normal CDF and each branch has its own forward and variance,
\begin{equation}
F_i = F\,e^{\,\mu_i+s_i^2/2},\qquad
B_i = \sigma_d^2 T + s_i^2,\qquad
d_{1i}=\frac{\log(F_i/K)+B_i/2}{\sqrt{B_i}},\quad d_{2i}=d_{1i}-\sqrt{B_i}.
\end{equation}
The branch variance $B_i$ is just the diffusion variance plus that scenario's jump
variance. Because $\sum_i w_iF_i=F$, put--call parity holds and, with no event, the
valuation collapses back to plain Black--Scholes. So instead of simulating, you
can split the world into $K$ branches at $t_e$, price each in closed form, and add them
up.

\section{Adding skew: the Heston model}

Jumps give the smile fat tails, but they don't produce the steady downward
\emph{skew} of equity options as out-of-the-money puts persistently pricier than
calls. To capture it, you can let volatility itself be random and
negatively correlated with the stock. This results in the standard Heston model, replacing the constant
$\sigma$ with a time-varying
variance $v_t$ that follows its own process,
\begin{equation}
dv_t = \kappa(\theta - v_t)\,dt + \xi\sqrt{v_t}\,dW_t^{v},
\qquad \mathrm{corr}(dW_t^S, dW_t^v)=\rho.
\end{equation}

\section{Combining them together}

Now, we have two indepdent pieces of randomness. We have both the Heston diffusion representing everyday moves,
as well asthe earnings jump modeling the one-off shock. 
A \emph{characteristic function}
(CF) is the Fourier transform of a return distribution, $\varphi(u)=\E[e^{iu X}]$, 
encoding the whole distribution in one function. It's useful here because Heston
has no closed-form density but does have a closed-form CF.

For independent R.V.s, the CF of the sum is the product of the CFs. So
the total return CF is just
\begin{equation}
\varphi_{\text{total}}(u)=\varphi_{\text{Heston}}(u)\cdot\varphi_{J}(u),
\qquad
\varphi_{J}(u)=\sum_{i} w_i\,e^{\,iu\mu_i-\frac12 u^2 s_i^2},
\end{equation}
where $\varphi_J$ is the CF of the Gaussian mixture. 

Combining Heston with a jump this way is a Bates model; the only difference here is
that the jump is the scheduled earnings mixture, applied only to expiries with
$t_e\le T$, rather than a Poisson jump.

\section{Pricing from the CF: COS}

With $\varphi_{\text{total}}$ in hand, prices come from the COS method
(Fang--Oosterlee): expand the risk-neutral density in a Fourier-cosine series on a
truncated range $[a,b]$ (set from the cumulants of $\varphi_{\text{total}}$). The
payoff's cosine coefficients are analytic, so the price collapses to a length-$N$
dot product against $\mathrm{Re}\{\varphi_{\text{total}}(u_k)e^{-iu_k a}\}$,
$N\approx 256$--$512$, $O(N)$ per strike. The engine is model-agnostic: switching
Black--Scholes $\to$ Heston $\to$ Bates only swaps the CF passed in
(\texttt{cos\_pricer.py}).

\section{Calibration}

Parameters are pinned sequentially from market structure rather than handed to one
optimizer, so the final fit only shapes the smile:
\begin{enumerate}\itemsep2pt
\item \textbf{Forward} from parity, $F=K+e^{rT}(C-P)$; invert Black-76 on the OTM
side for market IV per strike.
\item \textbf{Diffusion vol} is isolated by differencing two expiries that share
the event, which cancels it:
$\;\sigma_d^2=(\text{IV}_2^2 T_2-\text{IV}_1^2 T_1)/(T_2-T_1)$.
\item \textbf{Event variance} $V_e$ from the term-structure step: with
$\text{IV}(T)^2T\approx a+bT+\mathbf1_{\{t_e\le T\}}V_e$, fit $a,b$ on pre-event
expiries and read $V_e$ off a post-event one.
\item \textbf{Jump scale} matched to the ATM straddle \emph{price}, not its
variance --- a fat jump lifts the ATM straddle above its variance-implied level, so
matching variance leaves a uniform cross-strike bias that this removes.
\item \textbf{Remainder} (Heston $v_0,\xi,\rho$ and the mixture weights) fit to the
surface by minimizing IV-RMS. With the structural terms locked, only the skew and
smile are free.
\end{enumerate}

\section{Validation and overfitting}

The $0.00$ vol-point out-of-sample number is self-consistency, not forecasting: two
expiries bracketing the same event must carry the same forward event variance, so
recovering $V_e$ from one and matching the other only confirms forward-variance
additivity and a correct implementation.

The substantive test is against realized returns, which quotes can't referee.
Across 144 events (18 names, last 8 quarters), earnings days are $\sim$1.6\% of
trading days but a median 18\% of annual variance, and the reaction is
$\sim$3.5$\times$ a normal day --- the diffusion/event split is warranted. The
beat/miss assumption is not: on the pooled realized moves, BIC selects a single
fat-tailed component over 2 or 3. So the physical jump is best modeled unimodal;
the mixture remains useful degrees of freedom for fitting the \emph{implied} smile,
where asymmetry is genuinely priced, but it is not the return law.

\section{Scope}

This is not an alpha claim --- earnings and skew are already priced. The model gives a
consistent diffusion/skew/event decomposition of the surface, forecasts quantities
a single vol can't (e.g.\ the post-print IV crush), and flags where its marks
diverge from the market. Fits are in-sample; a systematic implied-vs-realized
premium study requires option history.

\end{document}
